\title{Group Sequence Policy Optimization}

\begin{abstract}
We present Group Sequence Policy Optimization (GSPO), a reinforcement learning algorithm designed for training large language models that emphasizes stability, efficiency, and high performance. Distinct from earlier approaches that rely on token-level importance ratios, GSPO constructs its importance ratio from sequence likelihoods and conducts all clipping, rewarding, and optimization at the sequence level. Our experiments reveal that GSPO outperforms the GRPO algorithm in both training speed and effectiveness, enhances the stability of RL training in Mixture-of-Experts (MoE) architectures, and offers the prospect of streamlining RL infrastructure. These advantages have played a pivotal role in the advancements seen in the most recent Qwen3 models.
\end{abstract}

\section{Introduction}
\label{sec:intro}

Reinforcement learning (RL) has become a central approach in the advancement of language models \citep{o1,dpsk-r1,qwq32b,qwen3}. Employing RL at large scale enables these models to enhance their abilities, allowing them to address complex tasks—including competition-grade mathematics and programming—by engaging in more extensive and intricate reasoning.

Effective scaling of RL with increased computational resources fundamentally depends on ensuring stable and resilient training processes. Despite advancements, leading RL algorithms such as GRPO \citep{grpo} are plagued by significant stability challenges when applied to large-scale language models, frequently causing total and irrecoverable model failure \citep{qwen3, minimax-m1}. Such instability significantly impedes attempts to further enhance the performance of language models via extended RL optimization.

In this study, we demonstrate that GRPO's instability originates from a core issue: the inappropriate use and resulting breakdown of importance sampling weights within its algorithmic structure. This flaw gives rise to training noise with substantial variance, which not only builds up as the response lengthens but is also exacerbated by the clipping procedure. These factors collectively lead to eventual model failure.

In response to these fundamental challenges, we introduce \textbf{Group Sequence Policy Optimization (GSPO)}, a novel reinforcement learning algorithm tailored for training large language models. GSPO's main advancement is its importance ratio, rigorously defined through sequence likelihood as established in prior work \citep{click}, and adheres to the foundational concepts of importance sampling. Furthermore, GSPO determines normalized rewards by evaluating the advantages among several responses to the same prompt, thus guaranteeing consistency between rewards assigned at the sequence level and the corresponding optimization objective.

Through our empirical analysis, we find that GSPO markedly outperforms GRPO in terms of training stability, efficiency, and overall effectiveness. Importantly, GSPO naturally addresses the stability issues that typically arise when applying RL to large Mixture-of-Experts (MoE) models, thereby obviating the necessity for intricate stabilization mechanisms and offering a pathway toward streamlined RL infrastructure. These advantages have played a crucial role in the substantial performance gains observed in the most recent Qwen3 models. We anticipate that GSPO will serve as a resilient and scalable algorithmic basis, fostering further progress in large-scale RL training for language models.

\section{Preliminaries}

\paragraph{Notation}
Throughout this work, we represent an autoregressive language model with parameters $\theta$ as the policy $\pi_\theta$. The symbol $x$ designates a query, while the set of all queries is indicated by $\mathcal{D}$. For any response $y$ generated in reply to $x$, the probability assigned by $\pi_\theta$ to $y$ given $x$ is expressed as $\pi_\theta (y | x)=\prod_{t=1}^{|y|} \pi_\theta (y_t | x, y_{<t} )$, where $|y|$ is the total token count in $y$. We employ a verifier $r$ to assess any query-response pair $(x, y)$, yielding a reward $r(x, y) \in [ 0, 1 ]$.

\paragraph{Proximal Policy Optimization (PPO)} PPO \citep{ppo} utilizes samples collected from the previous policy $\pi_{\theta_\text{old}}$, ensuring policy updates remain within a close neighborhood of the old policy by leveraging a clipping strategy. In particular, the method formulates policy optimization via the following objective (excluding the KL regularization term for simplicity, since it is not central to this discussion):

\begin{align}
\mathcal{J}_\text{PPO}(\theta) = \mathbb{E}_{ x \sim \mathcal{D},\, y \sim \pi_{\theta_\text{old}}( \cdot | x) }
\left[ \frac{1}{|y|} \sum_{t=1}^{|y|} 
\min \left( w_{t}(\theta) \widehat{A}_{t},  \, \mathrm{clip} \left( w_{t}(\theta), 1 - {\varepsilon}, 1 + {\varepsilon}\right) \widehat{A}_{t} \right)
\right],
\end{align}

Here, the token $y_t$’s importance ratio is given by $ w_{t}(\theta) = \frac{ \pi_{\theta} (y_{t} | x, y_{<t}) }{ \pi_{\theta_\text{old}} (y_{t} | x,y_{<t})} $. The advantage $\widehat{A}_{t}$ corresponding to $y_t$ is computed using a separate value model, while $\varepsilon$ specifies the allowable clipping interval for these importance ratios.

A principal practical limitation of PPO is its significant dependence on the value model. Notably, the value model often matches the policy model in scale, resulting in substantial demands on memory and computation. In addition, the overall performance of the algorithm is tightly linked to the accuracy of its value predictions. Given that achieving an accurate value model is already difficult, adapting such models to handle longer outputs and increasingly intricate tasks adds to these difficulties.

\paragraph{Group Relative Policy Optimization (GRPO)} GRPO \citep{grpo} eliminates the reliance on a value model by directly estimating the relative advantage among a collection of responses corresponding to an identical query. In particular, GRPO aims to maximize the following objective:

\begin{align}
\label{equ:grpo}
\mathcal{J}_\text{GRPO}(\theta) = \mathbb{E}_{ x \sim \mathcal{D},\, \{y_i\}_{i=1}^G \sim \pi_{\theta_\text{old}}( \cdot | x) }
\left[ \frac{1}{G} \sum_{i=1}^{G} \frac{1}{|y_i|} \sum_{t=1}^{|y_i|} 
\min \left( w_{i,t}(\theta) \widehat{A}_{i,t},  \, \mathrm{clip} \left( w_{i,t}(\theta), 1 - {\varepsilon}, 1 + {\varepsilon}\right) \widehat{A}_{i,t} \right)
\right],
\end{align}

where $G$ is the number of generated responses to each query $x$ (i.e., the group size), and the importance ratio $w_{i,t}(\theta)$ and advantage $\widehat{A}_{i,t}$ of token $y_{i,t}$ are:

\begin{align}
    w_{i,t}(\theta)=\frac{ \pi_{\theta} (y_{i,t} | x, y_{i,<t}) }{ \pi_{\theta_\text{old}} (y_{i,t} | x,y_{i,<t}) },\quad\ 
    \widehat{A}_{i,t} = \widehat{A}_{i} = \frac{ r(x, y_i) - \mathrm{mean} \left( \{ r(x, y_i) \}_{i=1}^G \right) }{ \mathrm{std} \left( \{ r(x, y_i) \}_{i=1}^G \right) },
\end{align}

Here, $w_{i,t}(\theta)$ denotes the importance weight calculated by taking the ratio of the probabilities assigned by the current policy, $\pi_\theta$, and the previous policy, $\pi_{\theta_\text{old}}$, for the token $y_{i,t}$, conditioned on $x$ and the preceding sequence $y_{i,<t}$. The advantage estimate, $\widehat{A}_{i,t} = \widehat{A}_{i}$, is obtained by normalizing the reward $r(x, y_i)$ for the $i$-th sample within a batch of size $G$: specifically, subtracting the mean reward over the batch and dividing by the standard deviation of the rewards, $\mathrm{std}\left(\{ r(x, y_i) \}_{i=1}^G \right)$.

respectively, where all the tokens in $y_i$ share the same advantage as $\widehat{A}_{i}$.

\section{Motivation}
\label{sec:motivation}

The increasing scale of models, the prevalence of sparsity (such as that found in Mixture-of-Experts architectures), and longer response sequences require utilizing substantial rollout batch sizes to fully leverage hardware during RL training. In order to enhance sample efficiency, it is common to subdivide these large batches into several smaller mini-batches for gradient computation. However, this process results in an off-policy training scenario: responses $y$ are drawn from a previous policy $\pi_{\theta_\text{old}}$ instead of the updated policy $\pi_\theta$ targeted for optimization. Consequently, algorithms such as PPO and GRPO employ clipping techniques to limit the influence of samples that are too far from the current policy, thereby ensuring stability in gradient updates.

Although techniques such as clipping are designed to address the off-policy gap, we highlight a deeper problem inherent to GRPO: \textit{the objective itself lacks well-posedness}. This challenge is especially pronounced during the training of large-scale models on tasks involving extended outputs, often resulting in severe model degradation. The root of this issue lies in the misuse of importance sampling weights within the GRPO framework. The foundational idea of importance sampling involves reweighting data gathered from a behavior policy $\pi_\text{beh}$ to approximate the expectation of a function $f$ with respect to a target policy $\pi_\text{tar}$:

\begin{align}
\mathbb{E}_{z \sim \pi_\text{tar}} \left[ f(z) \right] = \mathbb{E}_{z \sim \pi_\text{beh}} \left[ \frac{ \pi_\text{tar} (z) }{ \pi_\text{beh} (z)} \, f(z) \right].
\end{align}

Importantly, this approach depends on taking the average over a large number of samples ($N \gg 1$) drawn from the behavior distribution $\pi_\text{beh}$, enabling the importance weight $\frac{ \pi_\text{tar} (z) }{ \pi_\text{beh} (z)}$ to properly adjust for discrepancies between the distributions.

On the other hand, GRPO utilizes the importance weight $\frac{ \pi_{\theta} (y_{i,t} | x, y_{i,<t}) }{ \pi_{\theta_\text{old}} (y_{i,t} | x,y_{i,<t})}$ at each individual token index $t$. This weight is computed from a single observed token $y_{i,t}$, sampled from the next-token distribution $\pi_{\theta_\text{old}}( \cdot | x, y_{i, <t})$, and as such, does not adequately adjust for discrepancies between distributions as intended. Rather than facilitating accurate correction, it injects considerable high-variance noise into the gradient estimates, with this instability compounding throughout longer generated sequences and further intensified by the use of clipping. Empirical evidence indicates that these issues may cause the model to collapse in a manner that proves irreversible. Once collapse occurs, efforts to continue training—whether by restoring previous checkpoints, finely adjusting hyperparameters (like clipping thresholds), increasing sequence lengths, or altering RL signal queries—fail to recover model performance.

The observation outlined above highlights a central flaw in the design of GRPO. Specifically, the ineffectiveness of token-level importance weighting underscores a key concept: \textbf{the optimization objective must be aligned with the reward's granularity}. Given that rewards are assigned to whole sequences, implementing off-policy corrections on a per-token basis proves to be inappropriate. Consequently, we are led to abandon token-level objectives in favor of investigating approaches that employ importance weights and execute optimization at the \textit{sequence level} instead.

\section{Algorithm}

\subsection{GSPO: Group Sequence Policy Optimization}

Although the token-level importance weight $\frac{ \pi_{\theta} (y_{i,t} | x, y_{i,<t}) }{ \pi_{\theta_\text{old}} (y_{i,t} | x,y_{i,<t})}$ presents challenges within GRPO, our analysis indicates that, in language generation settings, the \textit{sequence-level} importance weight $\frac{ \pi_{\theta} (y | x) }{ \pi_{\theta_\text{old}} (y | x)}$ is theoretically well-founded. Specifically, it quantifies the extent to which a response $y$, drawn from $\pi_{\theta_\text{old}} (\cdot | x)$, diverges from $\pi_{\theta} (\cdot | x)$. This correspondence naturally connects to the sequence-level reward structure and provides an interpretable metric for the functioning of the clipping mechanism.

Based on this straightforward observation, we propose the \textbf{Group Sequence Policy Optimization (GSPO)} algorithm. GSPO employs the following sequence-level optimization objective:

\begin{align}
\mathcal{J}_\text{GSPO} (\theta) =
\mathbb{E}_{ x \sim \mathcal{D},\, \{y_i\}_{i=1}^G \sim \pi_{\theta_\text{old}}( \cdot | x) }
\left[ 
\frac{1}{G} \sum_{i=1}^{G}
\min \left( s_{i}(\theta)  \widehat{A}_{i},  \, \mathrm{clip} \left( s_{i}(\theta), 1 - {\varepsilon}, 1 + {\varepsilon} \right) \widehat{A}_{i} \right) 
\right],
\label{equ:gspo}
\end{align}

where we adopt the group-based advantage estimation:

\begin{align}
\widehat{A}_{i} = \frac{r(x, y_i) - \mathrm{mean} \left( \{ r(x, y_i) \}_{i=1}^G \right) }{ \mathrm{std} \left( \{ r(x, y_i) \}_{i=1}^G \right) },
\end{align}

and define the importance ratio $s_{i}(\theta)$ based on sequence likelihood \citep{click}:

\begin{align}
s_{i}(\theta) = \left( \frac{ \pi_{\theta} (y_i | x) }{ \pi_{\theta_\text{old}} (y_i | x)} \right)^{\frac{1}{|y_i|}}
=
\exp \left( \frac{1}{|y_i|} \sum_{t=1}^{|y_i|} \log \frac{ \pi_{\theta} (y_{i,t} | x, y_{i,<t}) }{ \pi_{\theta_\text{old}} (y_{i,t} | x,y_{i,<t})} \right).
\end{align}

Consequently, GSPO implements clipping at the response level, rather than on a token-by-token basis, to filter out samples that are excessively ``off-policy'' from contributing to the gradient estimation process. This approach ensures better alignment with the nature of both sequence-level rewards and the optimization objective. In addition, we use length normalization in $s_{i}(\theta)$ to mitigate variance and to maintain $s_{i}(\theta)$ within a consistent numerical scale. Without this normalization, substantial changes in the likelihood of a small number of tokens could lead to significant volatility in the sequence-level importance ratio, necessitating different clipping thresholds for responses of varying lengths. It is also worth mentioning that the clipping intervals in GSPO are generally of a different scale compared to those in prior methods such as GRPO, which is a result of the underlying differences in how the importance ratios are defined.

\subsection{Gradient Analysis}
\label{subsec:gradient}

We can derive the gradient of the GSPO objective as follows (clipping is omitted for brevity):

\begin{align}
\nabla_{\theta} \mathcal{J}_\text{GSPO} (\theta)
=&\ 
\nabla_{\theta} \mathbb{E}_{ x \sim \mathcal{D},\, \{y_i\}_{i=1}^G \sim \pi_{\theta_\text{old}}( \cdot | x) }
\left[ 
\frac{1}{G} \sum_{i=1}^{G} s_{i}(\theta) \widehat{A}_{i}
\right] \\
= &\ 
\mathbb{E}_{ x \sim \mathcal{D},\, \{y_i\}_{i=1}^G \sim \pi_{\theta_\text{old}}( \cdot | x) }
\left[ 
\frac{1}{G} \sum_{i=1}^{G}
s_{i}(\theta) \widehat{A}_{i}
\nabla_{\theta} \log s_{i}(\theta)
\right] \\
=&\ 
\mathbb{E}_{ x \sim \mathcal{D},\, \{y_i\}_{i=1}^G \sim \pi_{\theta_\text{old}}( \cdot | x) }
\left[ 
\frac{1}{G} \sum_{i=1}^{G} 
\left( \frac{ \pi_{\theta} (y_i | x) }{ \pi_{\theta_\text{old}} (y_i | x)} \right)^{\frac{1}{|y_i|}} \widehat{A}_{i} 
\frac{1}{|y_i|} \sum_{t=1}^{|y_i|} \nabla_{\theta} \log \pi_{\theta} (y_{i,t} | x, y_{i,<t}) 
\right].
\label{equ:gspo_grad}
\end{align}

For comparison, the gradient of the GRPO objective is as follows (note that $\widehat{A}_{i,t} = \widehat{A}_{i}$):

\begin{align}
\nabla_{\theta} \mathcal{J}_\text{GRPO} (\theta) 
=&\ 
\nabla_{\theta} \mathbb{E}_{ x \sim \mathcal{D},\, \{y_i\}_{i=1}^G \sim \pi_{\theta_\text{old}}( \cdot | x) }
\left[ \frac{1}{G} \sum_{i=1}^{G} \frac{1}{|y_i|} \sum_{t=1}^{|y_i|} 
w_{i,t}(\theta) \widehat{A}_{i,t}
\right] \\
=&\
\mathbb{E}_{ x \sim \mathcal{D},\, \{y_i\}_{i=1}^G \sim \pi_{\theta_\text{old}}( \cdot | x) }
\left[ \frac{1}{G} \sum_{i=1}^{G} \widehat{A}_{i} 
\cdot \frac{1}{|y_i|} \sum_{t=1}^{|y_i|} 
\frac{ \pi_{\theta} (y_{i,t} | x, y_{i,<t}) }{ \pi_{\theta_\text{old}} (y_{i,t} | x,y_{i,<t})} 
\nabla_{\theta} \log \pi_{\theta} (y_{i,t} | x, y_{i,<t})  
\right].
\end{align}

Accordingly, the primary difference between GSPO and GRPO is found in \textit{the manner in which they apply weights to the gradients of token log likelihoods}. GRPO utilizes a token-specific “importance weight,” namely $\frac{ \pi_{\theta} (y_{i,t} | x, y_{i,<t}) }{ \pi_{\theta_\text{old}} (y_{i,t} | x,y_{i,<t})}$, to modulate these gradients. The resulting weights—ranging across $(0, 1+\varepsilon]$ for $\widehat{A}_{i} >0$ and $[1-\varepsilon, +\infty)$ for $\widehat{A}_{i} < 0$—differ between tokens and may be significant; over the course of training, their cumulative effect can introduce unpredictability into the optimization process. In contrast, GSPO assigns identical weights to every token within a given response, thereby removing this potential source of instability present in GRPO.

\subsection{GSPO-token: A Token-level Objective Variant}

In contexts such as multi-turn RL, there may be a need for more granular advantage adjustments below the sequence level. Therefore, we propose a token-level adaptation of GSPO, referred to as \textbf{GSPO-token}, which facilitates the customization of advantages on a per-token basis:

\begin{align}
\mathcal{J}_\text{GSPO-token}(\theta)
=
\mathbb{E}_{ x \sim \mathcal{D},\, \{y_i\}_{i=1}^G \sim \pi_{\theta_\text{old}}( \cdot | x) }
\left[
\frac{1}{G} \sum_{i=1}^{G} \frac{1}{|y_i|} \sum_{t=1}^{|y_i|}
\min \left(
s_{i,t}(\theta) \widehat{A}_{i,t}, \ 
\mathrm{clip}\left( s_{i,t}(\theta), 1 - {\varepsilon}, 1 + {\varepsilon} \right) \widehat{A}_{i,t}
\right)
\right],
\label{equ:gspo-token}
\end{align}

where

\begin{align}
s_{i,t}(\theta) 
= \mathrm{sg} \left[ s_{i}(\theta) \right]  \cdot \frac{ \pi_{\theta} (y_{i,t} | x, y_{i,<t}) }{ \mathrm{sg} \left[ \pi_{\theta} (y_{i,t} | x, y_{i,<t}) \right] },
\end{align}

and $\mathrm{sg}[\cdot]$ denotes only taking the numerical value but stopping the gradient, corresponding to the \texttt{detach} operation in PyTorch. The gradient of GSPO-token can be derived as:

\begin{align}
\nabla_{\theta} \mathcal{J}_\text{GSPO-token}(\theta)
=&\
\nabla_{\theta} \mathbb{E}_{ x \sim \mathcal{D},\, \{y_i\}_{i=1}^G \sim \pi_{\theta_\text{old}}( \cdot | x) }
\left[
\frac{1}{G} \sum_{i=1}^{G}
\frac{1}{|y_i|} \sum_{t=1}^{|y_i|} 
s_{i,t}(\theta)\, \widehat{A}_{i,t}
\right] \\
=&\
\mathbb{E}_{ x \sim \mathcal{D},\, \{y_i\}_{i=1}^G \sim \pi_{\theta_\text{old}}( \cdot | x) }
\left[
\frac{1}{G} \sum_{i=1}^G s_i(\theta) \cdot \frac{1}{|y_i|} \sum_{t=1}^{|y_i|}
\widehat{A}_{i,t}\, \frac{\nabla_\theta \pi_\theta(y_{i,t} \mid x, y_{i,<t})}{\pi_\theta(y_{i,t} \mid x, y_{i,<t})}
\right] \\
=&\
\mathbb{E}_{ x \sim \mathcal{D},\, \{y_i\}_{i=1}^G \sim \pi_{\theta_\text{old}}( \cdot | x) }
\left[
\frac{1}{G} \sum_{i=1}^G \left( \frac{\pi_\theta(y_i \mid x)}{\pi_{\theta_\text{old}}(y_i \mid x)} \right)^{\frac{1}{|y_i|}}
\cdot \frac{1}{|y_i|} \sum_{t=1}^{|y_i|}
\widehat{A}_{i,t} \nabla_\theta \log \pi_\theta(y_{i,t} \mid x, y_{i,<t})
\right].
\label{equ:gspo-token_grad}
\end{align}

Observe that the expression $\frac{ \pi_{\theta} (y_{i,t} | x, y_{i,<t}) }{ \mathrm{sg} \left[ \pi_{\theta} (y_{i,t} | x, y_{i,<t}) \right] }$ evaluates to 1; as a result, $s_{i,t}(\theta)$ is, in terms of numerical value, the same as $s_{i}(\theta)$.

By juxtaposing Equation~\eqref{equ:gspo} with~\eqref{equ:gspo-token}, and Equation~\eqref{equ:gspo_grad} with~\eqref{equ:gspo-token_grad}, it becomes evident that GSPO-token and GSPO yield the same numerical results for the optimization objective, clipping mechanism, and theoretical gradient—provided that the advantage values assigned to each token in the sequence $y_i$ are set uniformly, namely, $\widehat{A}_{i,t} = \widehat{A}_{i}$ for all $t$. Nonetheless, GSPO-token permits greater adaptability by allowing the advantage to vary at the token level.

\section{Experiments and Discussion}

\subsection{Empirical Results}

We conduct experiments utilizing a cold-start model obtained by fine-tuning Qwen3-30B-A3B-Base, presenting both the reward progression during training and the model's performance trajectories on three benchmarks: AIME'24 (mean Pass@1 based on 32 samples), LiveCodeBench (202410-202502, mean Pass@1 across 8 samples), and CodeForces (using Elo Rating). For each training iteration in reinforcement learning, rollout batches are divided into four mini-batches to perform gradient updates. Within the GSPO framework, the lower and upper clipping thresholds in Equation~\eqref{equ:gspo} are set to 3e-4 and 4e-4, respectively. As a baseline, we utilize GRPO, specifying the left and right clipping intervals in Equation~\eqref{equ:grpo} as 0.2 and 0.27, values which have been carefully optimized for an equitable comparison. It is important to highlight that, for GRPO, the Routing Replay strategy is essential for stable MoE RL training—a topic further addressed in \S~\ref{subsec:routing_replay}—whereas \textit{GSPO eliminates this requirement}.

As illustrated in Figure~\ref{fig:results}, the training trajectory with GSPO remains consistently stable. Notably, \textbf{GSPO enables ongoing enhancements in model performance by scaling up training computation, periodically refreshing the query set, and prolonging the sequence length}. In addition, GSPO outperforms GRPO in terms of training efficiency, as it achieves higher accuracy and superior benchmark outcomes while utilizing equivalent computational resources and query counts. Finally, deploying GSPO for RL-based training of the latest Qwen3 models further substantiates its effectiveness, underscoring GSPO’s ability to maximize the scalability advantages of RL methods in large language model development.

\subsection{Curious Observation on Clipping Fractions}

An important difference between GSPO and GRPO lies in GSPO’s approach of clipping entire responses instead of performing clipping at the token level. Notably, as depicted in Figure~\ref{fig:clipping}, there is an observed gap of two orders of magnitude in the proportion of tokens that are clipped under GSPO as compared to GRPO (and modifying the clipping intervals does not affect this considerable difference). Interestingly, even though GSPO clips far more tokens—and therefore utilizes a smaller set for training or estimating gradients—it still attains better training efficiency than GRPO. This result, which may seem paradoxical given the substantially higher clipping rate, suggests that the token-based gradient estimations in GRPO tend to be noisy and suboptimal for leveraging training samples. On the other hand, GSPO’s reliance on sequence-level clipping produces a more robust and informative gradient signal for the learning process.

\subsection{Benefit of GSPO for MoE Training}
\label{subsec:routing_replay}

\paragraph{Background}
Unlike the reinforcement learning (RL) training of dense neural networks, the sparsely activated structure of Mixture of Experts (MoE) models brings forth specific challenges related to training stability. Notably, when implementing the GRPO algorithm, MoE models exhibit what we term as \textit{expert-activation volatility}, which hampers the stable convergence of RL optimization. Concretely, following one or more gradient updates, the subset of experts chosen to generate the same model output can shift considerably. To illustrate, in the case of the 48-layer Qwen3-30B-A3B-Base model, after each RL gradient step, for an identical rollout sample, approximately 10\% of the selected experts according to the updated policy $\pi_{\theta}$ differ from those selected under the previous policy $\pi_{\theta_\text{old}}$.
This characteristic, increasingly marked in deeper MoE architectures, leads to high variance in the token-level importance weights $w_{i,t}(\theta) = \frac{ \pi_{\theta} (y_{i,t} | x, y_{i,<t}) }{ \pi_{\theta_\text{old}} (y_{i,t} | x,y_{i,<t})}$, often rendering them unreliable, as further analyzed in \S~\ref{sec:motivation} and~\ref{subsec:gradient}. As a result, this volatility disrupts the standard RL training convergence process.

\paragraph{Our Previous Approach} In addressing this issue, our earlier method involved the use of the \textbf{Routing Replay} training paradigm. More precisely, we store the set of experts activated under $\pi_{\theta_\text{old}}$ and subsequently reuse these exact routing decisions in $\pi_{\theta}$ when evaluating the importance weights $w_{i,t}(\theta) = \frac{ \pi_{\theta} (y_{i,t} | x, y_{i,<t}) }{ \pi_{\theta_\text{old}} (y_{i,t} | x,y_{i,<t})}$. By adopting this strategy, both $\pi_{\theta} (y_{i,t} | x, y_{i,<t})$ and $\pi_{\theta_\text{old}} (y_{i,t} | x, y_{i,<t})$ utilize identical expert selections for each output token $y_{i,t}$, which reduces fluctuations in token-level importance ratios and helps guarantee stable training for the consistently activated parts of the network during gradient descent steps. As illustrated in Figure~\ref{fig:routing_replay}, Routing Replay proves vital for stable convergence in GRPO optimization applied to MoE architectures.

\paragraph{Benefit of GSPO} While Routing Replay assists GRPO training of MoE models in achieving convergence, its reliance on recycled routing modes increases both memory usage and communication demands, and may reduce the effective expressive power of the MoE model. By comparison, as depicted in Figure~\ref{fig:results}, GSPO obviates the need for Routing Replay, allowing conventional computation of importance ratios $s_i(\theta)$, leading to normal convergence and consistent optimization. The primary advantage of GSPO lies in its exclusive attention to the sequence likelihood, $\pi_{\theta} (y_i | x)$, rather than the granular likelihoods of individual tokens, $\pi_{\theta} (y_{i,t} | x, y_{i,<t})$. Given that the underlying language modeling ability of the MoE model is preserved throughout training, the sequence likelihood remains stable. In essence, GSPO directly addresses the instability of expert activation in MoE frameworks, thus removing the necessity for intricate methods such as Routing Replay. This results in a more streamlined and robust training workflow, enabling the MoE model to utilize its full capacity unencumbered by imposed limitations.

\subsection{Benefit of GSPO for RL Infrastructure}

Due to the differences in numerical precision between training frameworks (such as Megatron) and inference platforms (like SGLang and vLLM), practitioners often recalculate the likelihoods of generated responses under the prior policy $\pi_{\theta_\text{old}}$ using the training framework. However, GSPO relies exclusively on sequence-level likelihoods, rather than likelihoods computed at each token, for its optimization process. Sequence-level likelihoods are inherently more resilient to precision mismatches. As a result, GSPO enables the direct use of likelihood estimates provided by the inference engine during optimization, removing the need to recompute them via the training framework. This approach is particularly advantageous for applications including partial rollouts, multi-turn reinforcement learning, and workflows that separate training and inference operations.

\section{Conclusion}

We introduce Group Sequence Policy Optimization (GSPO), a novel reinforcement learning method designed for training large language models. Building on the foundation of importance sampling, GSPO establishes importance ratios derived from the likelihood of entire sequences and incorporates sequence-level operations such as clipping, rewarding, and optimization. In empirical comparisons, GSPO achieves markedly improved stability, efficiency, and overall performance relative to GRPO, with distinct advantages for large-scale reinforcement learning in MoE model scenarios. This approach has significantly contributed to the remarkable advancements seen in the latest Qwen3 models. With GSPO serving as a scalable algorithmic platform, our intention is to further extend reinforcement learning techniques, aiming for substantial breakthroughs in artificial intelligence.

\end{document}